\cleardoublepage                                  % 
\chapter{Antecedentes Y Estado del Arte}

\epigraph{“La matemática es la ciencia del orden y la medida, de bellas cadenas de razonamientos, todos sencillos y fáciles”.}{\textit{René Descartes}}

El panorama energético mundial atraviesa una transformación sin precedentes, impulsada por la necesidad de mitigar el cambio climático, reducir las emisiones de gases de efecto invernadero y transitar desde un modelo centralizado basado en combustibles fósiles hacia redes descentralizadas dominadas por fuentes de energía renovable (RES) \cite{muqeet2021energy}. La evolución tecnológica ha desplazado el enfoque desde arquitecturas centralizadas hacia modelos distribuidos y descentralizados. Los sistemas centralizados, aunque permiten una optimización global teórica, presentan riesgos de punto único de falla y limitaciones de escalabilidad debido a la carga computacional en un único procesador. Por el contrario, los enfoques descentralizados y distribuidos eliminan la dependencia de un controlador central, permitiendo que cada unidad tome decisiones basadas en información local, aunque esto puede comprometer la eficiencia global del sistema \cite{pamulapati2023review}.

En esta transición, las microrredes eléctricas emergen como una innovación crítica. Su capacidad para operar de forma autónoma o aislada (``\textit{off-grid}'') es fundamental para mejorar la resiliencia ante desastres naturales, asegurar el suministro en áreas rurales donde la extensión de la red principal no es técnica ni económicamente factible, y optimizar el consumo local de energía \cite{ibekwe2024microgrid}. Sin embargo, la gestión de una microrred en modo aislado presenta desafíos técnicos de alta complejidad. A diferencia de las redes conectadas a la utilidad principal, los sistemas \textit{off-grid} carecen de una fuente de referencia de tensión y frecuencia, lo que los hace vulnerables a las fluctuaciones de generación de recursos intermitentes como la energía solar fotovoltaica y eólica \cite{chae2014isolated}. 

En este contexto, y con el objetivo de garantizar la estabilidad del sistema y mantener el equilibrio entre generación y demanda, se han desarrollado diversas estrategias de gestión y control. Entre ellas se destacan el control jerárquico, los Sistemas Multiagente y los enfoques basados en la clasificación y priorización de cargas, los cuales permiten coordinar el comportamiento de los distintos componentes de la microrred frente a condiciones variables de operación \ref{sec:teoria_control}. 
Estas estragías a su vez implementan diferentes algoritmos de control, entre los cuales se destacan el control \textit{droop}, el basado en consenso, el control predictivo basado en modelo (MCP), controles adaptativos y técnicas basadas en aprendizaje automático.

Estás técnicas en conjunto a la adquisición de señales tanto de los generadores como de las cargas y la visualización forman el monitoreo de una microrred. Entre las tecnologías de monitoreo más comunes se encuentran:
\begin{itemize}
    \item \textbf{SCADA Sistemas de Supervisión, Control y Adquisición de Datos}: En microrredes, el SCADA se emplea para supervisar equipos como sensores, relés o \textit{breakers} y ejecutar órdenes de control, mejorando la confiabilidad y seguridad del sistema. Un SCADA puede encargarse de registrar medidas de corriente, voltaje o estado de interruptores, generar alarmas y permitir ajuste remoto, aunque con limitaciones de velocidad y complejidad creciente en redes grandes \cite{li2017_scada_microgrid}.
    \item \textbf{PMU (Unidades de Medición Fasorial)}: Son dispositivos que miden fasores de voltaje y corriente desde numerosas ubicaciones remotas de la red eléctrica a alta frecuencia con una fuente de sincronización común proporcionada por GPS. Las PMUs permiten un monitoreo en tiempo real de la estabilidad de la red con alta resolución temporal, lo que le otorga una ventaja sobre las mediciones SCADA tradicionales. El funcionamiento de la PMU se basa en el estándar IEEE C37.118 \cite{8975460_PMU_microgrid}.
    \item \textbf{Red de sensores IoT}: Plataformas basadas en Internet de las Cosas integran sensores distribuidos de corriente, voltaje u otros parámetros. Estos dispositivos inalámbricos de bajo costo pueden desplegarse masivamente, enviando datos en tiempo real a la nube o a controladores locales.  Las arquitecturas IoT pueden brindar flexibilidad y precisión en la gestión energética al combinar datos en tiempo real con modelos predictivos \cite{2025106158_iot_microgrid}.
    \item \textbf{Otras redes y protocolos}: Se emplean medidores inteligentes avanzados, unidades terminales remotas (RTU) y protocolos de IoT. Existen también desarrollos de monitoreo ``\textit{phygital}'' que combinan sensores físicos con agentes de software en cada nodo. En general, la estrategia de monitoreo debe garantizar alta observabilidad con latencia tolerable al nivel de control respectivo \cite{9023471_communication_microgrid}.
\end{itemize}

Para poder realizar la gestión, el control o monitoreo de una microrred es necesario la comunicación de los diferentes elementos que la componen. Dependiendo del nivel de control, los requisitos de comunicación pueden variar significativamente, por lo que es necesario seleccionar la tecnología de comunicación adecuada para cada caso. Se debe definir el medio de transmisión, el protocolo de comunicación y tener en cuenta la seguridad, los estándares y la latencia de la red para garantizar el funcionamiento deseado del sistema.
En medidos de transmisión se emplean redes cableadas (fibra óptica, par trenzado UTP, RS232, RS485, redes de comunicaciones por línea eléctrica) y redes inalámbricas (Wi-Fi, ZigBee, redes celulares 4G/5G). Cada tecnología tiene \textit{trade-offs}: la fibra ofrece baja latencia y alta fiabilidad pero costo alto; las redes Wi-Fi/ZigBee son económicas y escalables, pero sufren interferencia y riesgos de seguridad. Las redes celulares combinan baja latencia y cobertura amplia, aunque implican costo por licencias. Además debe se debe tener en cuenta el diseño de topologías redundantes como estrella o malla para evitar cortes. Mientras que a un nivel lógico se usan protocolos industriales como Modbus, DNP3, IEC 61850 (SAS para subestaciones y microrredes) y el estándar IEC 62351 para ciberseguridad en redes eléctricas. Muchos equipos modernos incorporan comunicaciones IP/TCP con cifrado TLS para interoperabilidad \cite{9023471_communication_microgrid}.